\chapter*{General Conclusion}
\label{chap:conclusion}
\addcontentsline{toc}{chapter}{General Conclusion}

In this thesis, we presented a hybrid digital twin for predictive maintenance of an FDM 3D printer. The system integrates a real‑time 3D visual model with a multi‑tier machine learning pipeline: rule‑based detection (Tier 0), Isolation Forest anomaly detection (Tier 1), and XGBoost classification (Tier 2).

We successfully built a low‑cost prototype (total hardware under 50 euros) using an MPU‑6050 vibration sensor, a KY‑013 temperature sensor, and an ESP8266 microcontroller. All processing runs on a local edge gateway (Raspberry Pi or laptop), with no cloud dependency. The dashboard is browser‑based and requires no installation, making it easy to use in small workshops, laboratories, and makerspaces.

Three case studies (nozzle clogging, motion‑axis vibrations, overheating) were carried out. The results show a weighted overall accuracy of 90.8\%, a 67\% reduction in print failure rate, and a 46\% decrease in unplanned downtime. The system also enables a shift from reactive to condition‑based maintenance.

\section*{Limitations and Future Work}
The current prototype has several limitations. It was tested on a single printer model (BCN3D+) with a limited set of failure scenarios. The dashboard handles only one printer at a time; a fleet management feature is planned but not yet implemented.

Future work includes:
\begin{itemize}
    \item Extending the system to support multiple printers simultaneously.
    \item Adding current sensing (ACS712) and vision‑based monitoring (camera) as additional sensor modalities.
    \item Developing a closed‑loop control that not only pauses the print but also adjusts parameters (e.g., nozzle temperature, fan speed) automatically.
    \item Conducting a long‑term field study to confirm economic benefits over several months.
\end{itemize}